% GENERATED FILE. Edit canonical lessons, then run npm run generate:books.
% canonical-source-hash: fnv1a64:1712bf2cd49dd5ca
% canonical-lessons: ZH-C232-chushi, ZH-C232-shoushang, ZH-C232-liuxue, ZH-C232-zhuang, ZH-C232-jiuhuche

\chapter{An Accident}
\label{ch:zh-an-accident}
\hlchaptermodality{232}

\begin{chapteropening}
\textbf{By the end of this chapter:} \emph{I can tell what happened in an accident: who got hurt, who was bleeding, what crashed and the ambulance.}

Complete the last lesson of chapter 232: I can tell what happened in an accident: who got hurt, who was bleeding, what crashed and the ambulance.
\end{chapteropening}

\section[\texorpdfstring{\zh{出事} (chūshì)}{出事 (chūshì)}]{\zh{出事} (chūshì) — to have an accident; for something to go wrong}

\label{lesson:ZH-C232-chushi}



\begin{quote}
Before the new one: say the Mandarin for a nurse, then the Mandarin for to prescribe.
\end{quote}

\subsection*{You'll want to know: \zh{出事}}
\textbf{\zh{出事}} — \emph{chūshì} — \textquotedblleft{}to have an accident; for something to go wrong\textquotedblright{}.

\begin{quote}
\textbf{\zh{前天雪很大这儿出事了}} \emph{qiántiān xuě hěn dà zhèr chūshì le} — The day before yesterday the snow was heavy, and there was an accident here.
\end{quote}

\begin{cousinweb}
\textbf{\zh{出事}} (\emph{chūshì}) is built from \zh{出} 'to come out' and \zh{事} 'a matter', both inherited from Middle Chinese, \emph{tsyhwit} and \emph{dzriH} (Guangyun, in Baxter's transcription), which glosses the second 'to employ' and 'to set up'; 'a matter' is the everyday sense, not the gloss.
\end{cousinweb}

\begin{grammarlens}[title={What you've built}]
One more word for talking about the past.
\end{grammarlens}

\subsection*{Your turn}
\begin{itemize}
\raggedright
  \item \emph{Say it:} \emph{chūshì}
  \item \emph{Say it:} \emph{chūshì}, once more
  \item \emph{Say it:} \emph{qiántiān xuě hěn dà zhèr chūshì le} — The day before yesterday the snow was heavy, and there was an accident here.
  \item \emph{Recall:} say the Mandarin for a nurse, then the Mandarin for to prescribe, then say \emph{chūshì} again
\end{itemize}

\begin{tcolorbox}[breakable,colback=teal!4,colframe=teal!35!black,title={Before you move on}]
What does \zh{出事} mean? (\textquotedblleft{}To have an accident; for something to go wrong\textquotedblright{}.) Say it once more.
\end{tcolorbox}
\section[shòushāng]{shòushāng — to get hurt; to be injured}

\label{lesson:ZH-C232-shoushang}



\begin{quote}
Before the new one: say the Mandarin for to prescribe, then the Mandarin for to have an accident.
\end{quote}

\subsection*{You'll want to know: shòushāng}
\emph{shòushāng} — \textquotedblleft{}to get hurt; to be injured\textquotedblright{}.

\begin{quote}
\emph{tā de tuǐ shòushāng le} — He hurt his leg.
\end{quote}

Its characters come later; for now it is a word for the ear.

\begin{cousinweb}
\emph{shòushāng} is built from \emph{shòu} 'to receive' and \emph{shāng} 'a wound', both inherited from Middle Chinese, \emph{dzyuwX} and \emph{syang} (Guangyun, in Baxter's transcription), which glosses the first 'to receive' and the second 'harm'.
\end{cousinweb}

\begin{grammarlens}[title={What you've built}]
One more word for talking about the past.
\end{grammarlens}

\subsection*{Your turn}
\begin{itemize}
\raggedright
  \item \emph{Say it:} \emph{shòushāng}
  \item \emph{Say it:} \emph{shòushāng}, once more
  \item \emph{Say it:} \emph{tā de tuǐ shòushāng le} — He hurt his leg.
  \item \emph{Recall:} say the Mandarin for to prescribe, then the Mandarin for to have an accident, then say \emph{shòushāng} again
\end{itemize}

\begin{tcolorbox}[breakable,colback=teal!4,colframe=teal!35!black,title={Before you move on}]
What does shòushāng mean? (\textquotedblleft{}To get hurt; to be injured\textquotedblright{}.) Say it once more.
\end{tcolorbox}
\section[liúxuè]{liúxuè — to bleed}

\label{lesson:ZH-C232-liuxue}



\begin{quote}
Before the new one: say the Mandarin for to have an accident, then the Mandarin for to get hurt.
\end{quote}

\subsection*{You'll want to know: liúxuè}
\emph{liúxuè} — \textquotedblleft{}to bleed\textquotedblright{}.

\begin{quote}
\emph{nǐ de shǒu liúxuè le} — Your hand's bleeding.
\end{quote}

Its characters come later; for now it is a word for the ear.

\begin{cousinweb}
\emph{liúxuè} is built from \emph{liú} 'to flow' and \emph{xuè} 'blood', both inherited from Middle Chinese, \emph{ljuw} and \emph{xwet} (Guangyun, in Baxter's transcription), which explains the first, after an older book, as water moving.
\end{cousinweb}

\begin{grammarlens}[title={What you've built}]
One more word for talking about the past.
\end{grammarlens}

\subsection*{Your turn}
\begin{itemize}
\raggedright
  \item \emph{Say it:} \emph{liúxuè}
  \item \emph{Say it:} \emph{liúxuè}, once more
  \item \emph{Say it:} \emph{nǐ de shǒu liúxuè le} — Your hand's bleeding.
  \item \emph{Recall:} say the Mandarin for to have an accident, then the Mandarin for to get hurt, then say \emph{liúxuè} again
\end{itemize}

\begin{tcolorbox}[breakable,colback=teal!4,colframe=teal!35!black,title={Before you move on}]
What does liúxuè mean? (\textquotedblleft{}To bleed\textquotedblright{}.) Say it once more.
\end{tcolorbox}
\section[zhuàng]{zhuàng — to crash into; to bump into}

\label{lesson:ZH-C232-zhuang}



\begin{quote}
Before the new one: say the Mandarin for to get hurt, then the Mandarin for to bleed.
\end{quote}

\subsection*{You'll want to know: zhuàng}
\emph{zhuàng} — \textquotedblleft{}to crash into; to bump into\textquotedblright{}.

\begin{quote}
\emph{tā qí chē zhuàngdàole yí ge rén} — He rode his bike into someone.
\end{quote}

Its characters come later; for now it is a word for the ear.

\begin{cousinweb}
\emph{zhuàng} is inherited from Middle Chinese \emph{draewngH} (Guangyun, in Baxter's transcription), glossed 'to strike a bell'; the dictionary glosses its other reading 'to knock against'.
\end{cousinweb}

\begin{grammarlens}[title={What you've built}]
One more word for talking about the past.
\end{grammarlens}

\subsection*{Your turn}
\begin{itemize}
\raggedright
  \item \emph{Say it:} \emph{zhuàng}
  \item \emph{Say it:} \emph{zhuàng}, once more
  \item \emph{Say it:} \emph{tā qí chē zhuàngdàole yí ge rén} — He rode his bike into someone.
  \item \emph{Recall:} say the Mandarin for to get hurt, then the Mandarin for to bleed, then say \emph{zhuàng} again
\end{itemize}

\begin{tcolorbox}[breakable,colback=teal!4,colframe=teal!35!black,title={Before you move on}]
What does zhuàng mean? (\textquotedblleft{}To crash into; to bump into\textquotedblright{}.) Say it once more.
\end{tcolorbox}
\section[jiùhùchē]{jiùhùchē — ambulance}

\label{lesson:ZH-C232-jiuhuche}



\begin{quote}
Before the new one: say the Mandarin for to bleed, then the Mandarin for to crash into.
\end{quote}

\subsection*{You'll want to know: jiùhùchē}
\emph{jiùhùchē} — \textquotedblleft{}ambulance\textquotedblright{}.

\begin{quote}
\emph{gǎnkuài jiào jiùhùchē} — Call an ambulance, quick!
\end{quote}

Its characters come later; for now it is a word for the ear.

\begin{cousinweb}
\emph{jiùhùchē} is built from \emph{jiù} 'to rescue', \emph{hù} 'to protect' and \zh{车} 'vehicle', all inherited from Middle Chinese, \emph{kjuwH}, \emph{huH} and \emph{tsyhae} (Guangyun, in Baxter's transcription), which glosses the first 'to protect' and 'to stop'.
\end{cousinweb}

\begin{grammarlens}[title={What you've built}]
One more word for talking about the past.
\end{grammarlens}

\subsection*{Your turn}
\begin{itemize}
\raggedright
  \item \emph{Say it:} \emph{jiùhùchē}
  \item \emph{Say it:} \emph{jiùhùchē}, once more
  \item \emph{Say it:} \emph{gǎnkuài jiào jiùhùchē} — Call an ambulance, quick!
  \item \emph{Recall:} say the Mandarin for to bleed, then the Mandarin for to crash into, then say \emph{jiùhùchē} again
\end{itemize}

\begin{tcolorbox}[breakable,colback=teal!4,colframe=teal!35!black,title={Before you move on}]
What does jiùhùchē mean? (\textquotedblleft{}Ambulance\textquotedblright{}.) Say it once more.
\end{tcolorbox}
